\documentclass[11pt]{article}

\usepackage[margin=1in]{geometry}
\usepackage[T1]{fontenc}
% Libertine text/math when available.
\IfFileExists{libertine.sty}{%
  \usepackage{libertine}%
  \IfFileExists{newtxmath.sty}{\usepackage[libertine]{newtxmath}}{}%
}{}
% Light monospace for inline code.
\IfFileExists{inconsolata.sty}{\usepackage[scaled=0.95]{inconsolata}}{%
  \IfFileExists{zi4.sty}{\usepackage[varqu,scaled=0.95]{zi4}}{}}
\usepackage{booktabs}
\usepackage{graphicx}
\usepackage{amsmath}
\usepackage{array}
\usepackage[table]{xcolor}
\usepackage[hidelinks]{hyperref}
\usepackage{xurl}
\usepackage{caption}

\captionsetup{font=small,labelfont=bf,skip=4pt}
\definecolor{hdrgray}{gray}{0.92}
\definecolor{keyrow}{RGB}{226,239,246}
\newcommand{\code}[1]{\texttt{#1}}
\newcommand{\best}[1]{\textbf{#1}}
\newcommand{\hdr}{\rowcolor{hdrgray}}
\newcommand{\key}{\rowcolor{keyrow}}
\setlength{\emergencystretch}{3em}

\newenvironment{ctab}{\par\medskip\centering\small
  \setlength{\tabcolsep}{5pt}\renewcommand{\arraystretch}{1.12}}{\par\medskip}

% Full-width abstract aligned with the body margins.
\renewenvironment{abstract}{%
  \par\medskip\noindent{\bfseries\abstractname.}\enspace\ignorespaces
}{\par\medskip}

\title{\textbf{Optimizing \code{get\_subtree} on MongoDB Community Edition:\\
A Measured Schema Study for the PageIndex Retrieval Workload}}
\author{Junyao Dong}
\date{}

\begin{document}
\maketitle

\begin{abstract}
On the PageIndex retrieval workload, \code{get\_subtree}---a materialized-path
range scan matching up to ${\sim}36$k nodes of a ten-million-node tree and
returning only small fields---is the one operation that stresses MongoDB. This
study measures which schema and index designs change it, all on Community
Edition. The gating cost is the \emph{count} of per-document fetches, not their
size: a non-covered scan reads and decodes every matched document regardless of
projection (MongoDB's product team confirmed the read path, and that newer
server versions do not change it), so shrinking documents buys nothing and
removing the fetch buys everything. The design that follows is full decoupling:
structure (\code{node\_id}, \code{path}, \code{parent\_id}, \code{depth}) in one
collection under a lean covering index, metadata (\code{title}, \code{summary})
in a key-value collection keyed by node id, text separate. The covered structure
scan runs with zero document fetches at $15.5$\,ms P50 / $168$\,ms P95 in the
ten-million-node MongoDB-only ablation. A separate symmetric three-million-node
matrix gives PostgreSQL the equivalent covering index and narrow structure
table: MongoDB remains $4.35{\times}$/$3.91{\times}$ slower at P50/P95 in the
covered arm and $4.65{\times}$/$4.07{\times}$ slower in the narrow-structure
arm. The study therefore demonstrates a MongoDB optimization, not cross-engine
parity. Structural denormalization may reduce the remaining per-key cost but is
unmeasured.
\end{abstract}

\section{The problem}

The PageIndex retriever reaches storage through four calls: a point lookup, an
expand-children, a content fetch, and \code{get\_subtree}, which pulls a node's
descendants to render the tree view and hand a block to the language model. On
the ten-million-node dataset the first three are sub-millisecond on MongoDB and
need no work. \code{get\_subtree} is the exception: a large subtree matches
about 36{,}000 nodes, and returning tens of thousands of records exercises the
engine's per-record read cost like nothing else in the workload. On a monolithic
node document---structure, \code{title}, \code{summary}, and leaf \code{text} in
one document per node---the scan performs one \code{FETCH} per match. This study
asks which schema removes avoidable work inside MongoDB. Cross-engine claims are
kept separate and use only the symmetric matrix in Table~\ref{tab:fair}; no
optimized MongoDB result is divided by an unoptimized PostgreSQL baseline.

\paragraph{Conditions.} The MongoDB schema/index ablation uses the
ten-million-node tree: one measuring client, the collection freshly loaded and
fully resident in the WiredTiger cache, and the same 200 seeded subtree paths
(depth ${\le}3$, fixed seed) across MongoDB variants. Those results are directly
comparable only within that ablation. The cross-engine matrix is a separate
campaign on a clean 16-vCPU/30-GiB host, a three-million-node tree, one engine
per invocation, and the same 200 paths and returned ids for MongoDB and
PostgreSQL. The tested tail is not caused by interference: driving
the scan while 16--64 background processes hammer point lookups at the same
\code{mongod} moves P95 by only ${\sim}1.1{\times}$, with zero application-thread
evictions and the collection (${\sim}17$\,GB) fully resident throughout.

\begin{table}[h]
\caption{Symmetric ids-only \code{get\_subtree} matrix, 3M nodes, average
8{,}840.4 returned ids (P50 / P95 ms). Every row gives both engines the same
layout and index capability; the ratio is MongoDB / PostgreSQL.}
\label{tab:fair}
\begin{ctab}
\begin{tabular}{@{}lrrr@{}}
\toprule
\hdr Symmetric arm & MongoDB & PostgreSQL & Ratio \\
\midrule
Monolith + plain \code{path} index & 2.391 / 19.201 & \best{0.648 / 6.817} & 3.69 / 2.82$\times$ \\
\key Monolith + covering index & 1.291 / 10.692 & \best{0.297 / 2.732} & 4.35 / 3.91$\times$ \\
Narrow structure + covering index & 1.269 / 10.616 & \best{0.273 / 2.610} & 4.65 / 4.07$\times$ \\
\bottomrule
\end{tabular}
\end{ctab}
\end{table}

This server-to-server table is the report's cross-engine comparison. SQLite has
equivalent symmetric arms but is in-process, so it is not used for the primary
ratio. DuckDB has no independent covered or narrow arm: columnar projection
already reads only \code{node\_id}, making covered equal to naive by
construction.

\paragraph{PostgreSQL hierarchy-layout sensitivity.}
Table~\ref{tab:fair} deliberately holds the materialized-path representation
fixed; it is a matched comparison, not a best-schema tournament. The label
``PostgreSQL (JSONB)'' is also easy to misread here: \code{path} and
\code{node\_id} are typed SQL columns, and the covered subtree query never
reads the \code{JSONB} payload. We nevertheless ran a separate PostgreSQL-only
diagnostic starting from the real API's root \code{node\_id}, while retaining
the microbenchmark's unbounded descendant-ID output. It compares five
representations on PostgreSQL 16.14, the same three-million-node tree, and 200
roots, with three cyclically interleaved repeats per root.

\begin{table}[t]
\caption{PostgreSQL-only hierarchy-layout sensitivity, P50/P95 ms. This table
is not used for a MongoDB/PostgreSQL ratio.}
\label{tab:pg-hierarchy}
\begin{ctab}
\begin{tabular}{@{}lrl@{}}
\toprule
\hdr Layout & P50 / P95 & Measured access path \\
\midrule
Typed materialized path & 0.401 / 4.644 & two covered B-tree probes \\
DFS preorder interval & 0.365 / \best{4.323} & two covered B-tree probes \\
Closure table & \best{0.358} / 4.515 & one covered B-tree probe \\
PostgreSQL \code{ltree} & 0.655 / 6.836 & GiST bitmap + heap scan \\
Recursive adjacency & 1.602 / 21.132 & recursive union + child probes \\
\bottomrule
\end{tabular}
\end{ctab}
\end{table}

All five layouts returned identical IDs for all 200 roots. The result is more
limited than the phrase ``PostgreSQL-native layout'' suggests. The current
typed path is already within 9\% at P50 and 7\% at P95 of preorder; the latter
is a generic read-optimized nested-set encoding, not a PostgreSQL-specific
type. PostgreSQL's supplied hierarchy type, \code{ltree}, was slower here
because its descendant operator used a GiST bitmap scan and revisited heap
pages instead of producing the IDs through a zero-heap-fetch B-tree scan.
This arm used a 100-byte GiST signature, rather than the eight-byte default, to
make the signature filter more precise; the default was slower in the medium
control.
Closure removes the root-to-range lookup but expands 3 million nodes into
17.1 million ancestor/descendant rows (5.71 rows per node) and still does not
improve P95 over preorder. Simple integer preorder is therefore a reasonable
optional sensitivity for a bulk-built, mostly static tree, not a missing
baseline that invalidates Table~\ref{tab:fair}; inserts can require relabeling
or a gap-labeling scheme. None of these PostgreSQL-only values is divided by a
MongoDB value.

\section{Root cause: the count of fetches, not their size}

A secondary index in MongoDB (here on \code{path}) is a separate B-tree holding
only the indexed value and a record id. To return any other field, WiredTiger
must \emph{fetch} the full BSON document and apply the projection to it. A point
lookup is one such fetch and is cheap; \code{get\_subtree} is tens of thousands
of them. MongoDB's product team confirmed this read path when we raised it:
without a covering index, the server reads and decodes each whole matched
document regardless of how small the projection is, and newer server versions do
not change it (the 8.0 SBE included). The data is resident and WiredTiger caches
pages uncompressed, so the cost is neither disk I/O nor decompression---it is
per-document work, times subtree size.

Two measurements pin the cost to the count of fetches rather than their size.
Splitting the wide \code{text} field into its own collection (the Subset
Pattern) shrinks every scanned document by ${\sim}75\%$
($4.97\to1.25$\,GB of data) and leaves the scan unchanged ($419\to419$\,ms P95
for the view projection). Going further to ${\sim}150$-byte topology-only
documents---an order of magnitude smaller again---still pays most of the fetch
cost: $260$\,ms P95 (Table~\ref{tab:kv}) against the monolithic layouts'
$284$--$301$\,ms. Locating and
materializing ${\sim}36$k documents dominates regardless of what they weigh, so
only removing the fetch helps. Two corollaries follow. The split is still
correct---\code{get\_subtree} never needs \code{text}, and on a
cache-constrained host, unlike this one, inline text would push the hot working
set out of RAM---it is just not a latency fix. And the relational engines win
this one scan for the same reason: a row store materializes only the requested
columns and a column store reads only the requested column, so both skip
per-node work MongoDB pays.

\section{Why one document per node cannot deliver the scan}\label{sec:mono}

Keeping the monolithic layout and adding indexes was measured first, and the
options bracket the tension. A lean \code{\{path,node\_id\}} covering index
answers the id-only scan index-only (\code{explain}: \code{PROJECTION\_COVERED},
\code{totalDocsExamined:\,0}) at $191$\,ms P95 versus the MongoDB-only
non-covered result of $301$\,ms---this is the shape MongoDB's product team recommended,
the covering index rather than the split, with the write-once tree paying the
index build once at ingest. But the tree view needs \code{title} and
\code{summary}. Covering them too (\code{\{path,node\_id,title,summary\}}) still
scans covered and beats the non-covered view ($419\to350$\,ms), yet the
multi-kilobyte \code{summary} key inflates the index to $4.66$\,GB---the size of
the collection's own data---because prefix compression deduplicates the shared
\code{path} prefixes but does nothing for fat values. Resolving the view with a
covered id scan plus a batched \code{\$in} back into the same collection is
worse than not covering at all ($1{,}314$\,ms P95) in this harness: its fixed
1{,}000-id chunk size creates ${\sim}36$ round trips that re-fetch the same fat
documents. The current short ids would form only about a 0.62\,MiB BSON command,
so this chunk count is not imposed by MongoDB's 16\,MiB command limit. One
layout therefore cannot hold both a lean index and the view fields; the candidate
schema stops trying.

\section{Candidate schema: structure/metadata decoupling,
measured}\label{sec:kv}

Structure (\code{node\_id}, \code{path}, \code{parent\_id}, \code{depth}) lives
in a topology collection of ${\sim}150$-byte documents under the lean
\code{\{path,node\_id\}} covering index; metadata (\code{title}, \code{summary})
lives in a key-value collection keyed by \code{(tree\_id,node\_id)} (or an
equivalently namespaced \code{\_id}); \code{text} stays in its own collection.
\code{get\_subtree} becomes a covered scan over the
structure collection, and the view's metadata resolves by \code{\_id} point
reads for the nodes actually rendered.

\begin{table}[h]
\caption{\code{get\_subtree} on the decoupled schema, 10M tree (one client, ms),
same 200 subtree paths. Anchor rows re-measure the no-\code{text} monolithic
layout in the same run. The resolution row prices the worst case: metadata for
\emph{every} node of a ${\sim}36$k-node subtree in one call.}
\label{tab:kv}
\begin{ctab}
\begin{tabular}{@{}p{7.0cm}rrl@{}}
\toprule
\hdr Read & P50 & P95 & Note \\
\midrule
Anchor: no-\code{text} monolith, view (\code{FETCH}) & 38.3 & 421 & \\
Anchor: no-\code{text} monolith, id (covered) & 17.3 & 181 & \\
Structure scan, id, no covering index & 25.2 & 260 & \code{FETCH} per node \\
\key Structure scan, id (covered) & 15.5 & 168 & \code{totalDocsExamined:\,0} \\
${+}$ metadata for all nodes, 1k-id chunks & 96.3 & 1034 & harness: ${\sim}36$ calls \\
% TODO(kv-run): add the server-side $lookup and clustered-metadata rows when the
% remaining bench_subset_kv.py phases have run (kv_lookup, kv_in_clu, kv_lookup_clu).
\bottomrule
\end{tabular}
\end{ctab}
\end{table}

The covered structure scan is the fastest \code{get\_subtree} measured in this
MongoDB-only 10M study---$15.5$\,ms P50 / $168$\,ms P95, with zero document
fetches. It is not a cross-engine result; Table~\ref{tab:fair} is the symmetric
comparison. The split costs the fast path little because the two traps MongoDB's
product team flagged are avoidable by
construction: the index carries exactly the projected fields (a missing field,
or \code{\_id} left in the projection, silently reinstates the \code{FETCH}),
and the index build is deferred to after bulk load, so the write-once tree pays
its maintenance once. Bulk metadata resolution stays the worst case, not the hot
path: \code{title}+\code{summary} for all ${\sim}36$k nodes via chunked
\code{\$in} against the key-value store costs ${\sim}1$\,s at P95---better than
the same join against the monolithic collection, but still ${\sim}36$k point
lookups behind the harness's ${\sim}36$ calls. This is a stress row rather than a
production endpoint measurement; the actual formatter's fields, result bounds,
and batching still need to be measured.

On disk the navigation working set is ${\sim}1.7$\,GB---$0.28$\,GB structure
data, $0.57$\,GB covering index (prefix compression deduplicates the shared
\code{path} prefixes), $0.69$\,GB metadata, $0.14$\,GB \code{\_id}
index---against $5.3$\,GB for the monolithic collection with \code{text}
inline. At production scale, with many trees sharing a host, that difference is
what keeps navigation cache-resident.

\section{Further structural denormalization (untested)}

The covered scan still performs per-key index work. Three designs may reduce it,
all viable because the tree is largely write-once,
so their cost is paid at load. \emph{Nested Sets}: number nodes
\code{left}/\code{right} in a DFS at ingest, making a subtree one range query;
MongoDB recommends it for static trees, and its $O(n)$ renumbering on insert
fits a load-then-serve deployment or batched reindexing, not steady growth.
\emph{Tree-order \code{\_id} on a clustered collection}: documents stored
physically in tree order turn the subtree range into a sequential read (the
optimizer needs a hint when a secondary index exists). \emph{Precomputed subtree
documents}: materialize each bounded subtree's structure+summary view at ingest,
turning \code{get\_subtree} into a point read; buckets of a few levels stay
under the 16\,MiB document limit, at the cost of storage amplification across
overlapping subtrees. They are unmeasured, so no claim that they close the fair
cross-engine gap follows. Whether that complexity is worthwhile is a product
decision based on a complete endpoint measurement. (\code{\$graphLookup} is not a candidate: it re-traverses
\code{parent\_id}, still materializes documents, and in 7.0 is capped at
100\,MB, ignoring \code{allowDiskUse}.)

\section{What a version bump or a license would not change}

MongoDB's product team confirmed the non-covered read path is unchanged in
newer server versions, and nothing in 8.0--8.3 targets this shape: the Express
path accelerates point/equality lookups only, block processing is scoped to
time-series collections, and the \code{\$graphLookup} disk-spill fix
(SERVER-23980, 8.2.0) helps only the recursive alternative that loses to the
range scan anyway. On Atlas, a column-store index would read only projected
columns, but PageIndex never scans the whole corpus; Search and Vector Search
serve relevance retrieval, not a range projection. Every lever this workload
needs---schema, index, storage configuration---is on Community.

\section{What remains open}\label{sec:open}

Three items. The remaining bulk-resolution variants---a server-side
\code{\$lookup} join from the structure scan into the key-value store (one
round trip instead of ${\sim}36$), and a clustered organization of the metadata
collection---are queued in the same harness. The Subset split under a
\emph{cache-constrained} host is the most useful retest: this study ran fully
resident, which by construction cannot show a residency benefit, so a run with
\code{--wiredTigerCacheSizeGB} below the working set would price the split's
real production value. And the structural denormalizations above remain
unpriced. The practical conclusion is limited: $168$\,ms P95 is a useful 10M
MongoDB-only zero-fetch capacity point, while the symmetric 3M matrix still
shows a roughly $4{\times}$ server-to-server gap for the ids-only scan.

\section{Conclusion}

Part of MongoDB's \code{get\_subtree} cost is layout-induced, and the
measurements show how to remove that avoidable part. The cost is the count of per-document fetches---confirmed by
MongoDB's product team, unchanged by newer versions, indifferent to document
size---so the fix is to take documents off the scan path entirely: structure
under a lean covering index (zero fetches, $15.5$\,ms P50 / $168$\,ms P95 on
ten million nodes in the MongoDB-only ablation), metadata as tenant-qualified
key-value reads, and text fetched only for selected nodes. The fair 3M matrix
does not show parity: with equivalent covering indexes MongoDB remains
$4.35{\times}$/$3.91{\times}$ slower at P50/P95, and with symmetric narrow
structure stores it remains $4.65{\times}$/$4.07{\times}$ slower. What remains
is the cache-constrained retest, complete endpoint measurement, and the queued
resolution variants.

\section*{References}
\small
\begin{enumerate}\itemsep2pt
\item Subset Pattern. \url{https://www.mongodb.com/docs/manual/data-modeling/design-patterns/group-data/subset-pattern/}
\item Model tree structures with Nested Sets. \url{https://www.mongodb.com/docs/v7.0/tutorial/model-tree-structures-with-nested-sets/}
\item Model tree structures with Materialized Paths. \url{https://www.mongodb.com/docs/manual/tutorial/model-tree-structures-with-materialized-paths/}
\item Clustered collections. \url{https://www.mongodb.com/docs/v7.0/core/clustered-collections/}
\item WiredTiger storage engine (cache holds data uncompressed; cache sizing). \url{https://www.mongodb.com/docs/manual/core/wiredtiger/}
\item Covered queries; explain results (covered = IXSCAN with no FETCH). \url{https://www.mongodb.com/docs/manual/reference/explain-results/}
\item \code{\$graphLookup} memory limit. \url{https://www.mongodb.com/docs/v7.0/reference/operator/aggregation/graphlookup/}
\item MongoDB 8.0 release notes (Express path; block processing). \url{https://www.mongodb.com/docs/manual/release-notes/8.0/}
\item SERVER-23980 (\code{\$graphLookup} disk spill, fixed 8.2.0); SERVER-102453 (reverted from 8.1). \url{https://jira.mongodb.org/browse/SERVER-23980}
\item MongoDB Community Edition feature scope (column-store and analytics are Atlas/Enterprise). \url{https://www.mongodb.com/products/self-managed/community-edition}
\end{enumerate}

\end{document}
